\documentclass[exercice,a4]{thedocument}%
\usepackage{pstricks}%
\usepackage{multido}%
\usepackage{tabularx}%
\usepackage[np]{numprint}%
\usepackage{enumitem}%
\usepackage{colortbl}%
\startdatakeys
    \source{25GENMATAA1}
    \BO{2026}
    \cycle{4}
    \competencesDuSocle{CA,RA}
    \objectifsApprentissage{C4A1aa,C4A2ad,C4B2cb,C4D2aa,C4D3ac}
    \motsCles{périmètre,expression littérale,tableur,équation}
\enddatakeys
\begin{document}%
    

Dans cet exercice, toutes les longueurs sont exprimées en centimètres.

On considère :
\begin{itemize}[label=$\bullet$]
\item le rectangle ABCD tel que AD = $x$ et AB = $16 - 2x$ ;
\item le carré EFGH tel que EF $= 2x$.
\end{itemize}

\psset{unit=0.85cm,arrowsize=2pt 3}
\begin{center}
\begin{pspicture}(0,-1)(15,5)
%\psgrid
\psframe(0,2)(10,4)
\uput[d](0,2){A}\uput[u](0,4){D}\uput[d](10,2){B}\uput[u](10,4){C}
\psline[linewidth=0.8pt]{<->}(-0.5,2)(-0.5,4) \uput[l](-0.5,3){$x$}
\psline[linewidth=0.8pt]{<->}(0,1)(10,1) \uput[d](5,1){$16 - 2x$}
\psframe(0,2)(0.2,2.2)\psframe(0,4)(0.2,3.8)
\psframe(10,4)(9.8,3.8) \psframe(10,2)(9.8,2.2)
{\psset{linecolor=blue}
\multido{\n=4.9+0.1}{3}{\psline(\n,1.8)(\n,2.2)}
\multido{\n=4.9+0.1}{3}{\psline(\n,3.8)(\n,4.2)}
\multido{\n=2.9+0.1,\na=3.1+0.1}{2}{\psline(-0.2,\n)(0.2,\na)}
\multido{\n=2.9+0.1,\na=3.1+0.1}{2}{\psline(9.8,\n)(10.2,\na)}
}
\psframe(11,0)(15,4)
\psline[linewidth=0.8pt]{<->}(11,-0.8)(15,-0.8) \uput[d](13,-0.8){$2x$}
\uput[d](11,0){E}\uput[u](11,4){H}\uput[d](15,0){F}\uput[u](15,4){G}
\psframe(11,0)(11.2,0.2) \psframe(15,0)(14.8,0.2)
\psframe(15,4)(14.8,3.8) \psframe(11,4)(11.2,3.8)
{\psset{linecolor=blue}
\multido{\n=0.9+0.1,\na=1.1+0.1}{2}{\psline(10.8,\n)(11.2,\na)}
\multido{\n=2.9+0.1,\na=3.1+0.1}{2}{\psline(10.8,\n)(11.2,\na)}
\multido{\n=0.9+0.1,\na=1.1+0.1}{2}{\psline(14.8,\n)(15.2,\na)}
\multido{\n=2.9+0.1,\na=3.1+0.1}{2}{\psline(14.8,\n)(15.2,\na)}
\multido{\n=-0.15+0.10,\na=0.05+0.10}{2}{\psline(11.8,\n)(12.2,\na)}
\multido{\n=-0.15+0.10,\na=0.05+0.10}{2}{\psline(13.8,\n)(14.2,\na)}
\multido{\n=3.85+0.10,\na=4.05+0.10}{2}{\psline(11.8,\n)(12.2,\na)}
\multido{\n=3.85+0.10,\na=4.05+0.10}{2}{\psline(13.8,\n)(14.2,\na)}}
\psline(13,-0.2)(13,0.2) \psline(13,3.8)(13,4.2) 
\psline(10.8,2)(11.2,2) \psline(14.8,2)(15.2,2) 
\end{pspicture}
\end{center}

\bigskip

\textbf{PARTIE A :} Dans cette partie, $x = \np[cm]{1,5}$.

\medskip

\begin{enumerate}
\item Calculer le périmètre du carré EFGH.
\item Calculer AB.
\item Construire en vraie grandeur le rectangle ABCD.
\item Les périmètres de ABCD et EFGH sont-ils égaux ?
\end{enumerate}

\bigskip

\textbf{PARTIE B :} Dans cette partie, on cherche pour quelle(s) valeur(s) de $x$, le périmètre du rectangle est égal au périmètre du carré.

\medskip

\begin{enumerate}
\item Pour essayer de répondre au problème, on utilise la feuille de calcul suivante:

\smallskip

\begin{center}
\renewcommand \arraystretch {1}
\begin{tabularx}{\linewidth}{|>{\cellcolor{lightgray}} c|m{3.85cm}|*{6}{>{\centering \arraybackslash}X|}}
\cline{2-8}
\rowcolor[gray]{0.7}\multicolumn{1}{c|}{\cellcolor{white}}&\centering A&B&C&D&E&F&G\\\hline
1&Valeur de $x$&1&2&3&4&5&6\\\hline
2&Périmètre du carré&8&16&24&32&40&48\\\hline
3&Périmètre du rectangle&30&28&26&24&22&20\\\hline
\end{tabularx}
\end{center}

\begin{enumerate}
\item Quel formule a-t-on pu saisir dans la cellule B2 avant de l'étirer jusqu'à G2?
\item Ce tableau nous permet-il de trouver une valeur de $x$ pour laquelle les deux périmètres sont égaux ?
\end{enumerate}

\item 
	\begin{enumerate}
		\item Montrer que le périmètre du rectangle peut s'écrire $-2x + 32$.
		\item Déterminer la solution au problème par la résolution d'une équation.
	\end{enumerate}
\end{enumerate}

\newpage
\showthesource
\end{document}
